% =====================================================================
%  Anatomía de un LLM — Fase I — Módulo 1
%  Modelos de lenguaje probabilísticos
%  Compilar: latexmk -xelatex M01-modelos-de-lenguaje.tex
% =====================================================================
\documentclass[11pt,a4paper,twoside,openany]{book}
\usepackage{../estilo-apuntes}

\newcommand{\ApuntesCurso}{Anatomía de un LLM}
\newcommand{\ApuntesModulo}{Módulo 1}

\hypersetup{pdftitle={Módulo 1. Modelos de lenguaje probabilísticos (resumen)},pdfauthor={Anatomía de un LLM}}

\makeatletter
\newlabel{ej:luna}{{1.5}{2}{Tokenización por caracteres}{ejemplo.5}{}}
\newlabel{ej:luna@cref}{{[ejemplo][5][]1.5}{[1][1][]2}}
\newlabel{eje:1-contar}{{1.1}{2}{}{ejercicio.1.1}{}}
\newlabel{eje:1-contar@cref}{{[ejercicio][1][1]1.1}{[1][2][]2}}
\newlabel{eje:1-inyectiva}{{1.2}{2}{}{ejercicio.1.2}{}}
\newlabel{eje:1-inyectiva@cref}{{[ejercicio][2][1]1.2}{[1][2][]2}}
\newlabel{eje:1-granularidad}{{1.3}{3}{}{ejercicio.1.3}{}}
\newlabel{eje:1-granularidad@cref}{{[ejercicio][3][1]1.3}{[1][2][]3}}
\newlabel{ej:bigramas}{{2.7}{5}{Un modelo de bigramas}{ejemplo.7}{}}
\newlabel{ej:bigramas@cref}{{[ejemplo][7][]2.7}{[1][5][]5}}
\newlabel{fig:cadena}{{2.1}{6}{Diagrama de transición del modelo de bigramas del \cref {ej:bigramas}. Cada arco $i\to j$ lleva $P(j\mid i)$; las probabilidades que salen de cada estado suman~1}{figure.2.1}{}}
\newlabel{fig:cadena@cref}{{[figure][1][2]2.1}{[1][5][]6}}
\newlabel{eje:2-casillas}{{2.1}{6}{}{ejercicio.2.1}{}}
\newlabel{eje:2-casillas@cref}{{[ejercicio][1][2]2.1}{[1][6][]6}}
\newlabel{eje:2-vf}{{2.2}{6}{}{ejercicio.2.2}{}}
\newlabel{eje:2-vf@cref}{{[ejercicio][2][2]2.2}{[1][6][]6}}
\newlabel{eje:2-sol-matriz}{{2.3}{6}{}{ejercicio.2.3}{}}
\newlabel{eje:2-sol-matriz@cref}{{[ejercicio][3][2]2.3}{[1][6][]6}}
\newlabel{eje:2-sol-prob}{{2.4}{7}{}{ejercicio.2.4}{}}
\newlabel{eje:2-sol-prob@cref}{{[ejercicio][4][2]2.4}{[1][6][]7}}
\newlabel{eje:2-longitud}{{2.5}{7}{}{ejercicio.2.5}{}}
\newlabel{eje:2-longitud@cref}{{[ejercicio][5][2]2.5}{[1][7][]7}}
\newlabel{eje:2-sin-fin}{{2.6}{7}{}{ejercicio.2.6}{}}
\newlabel{eje:2-sin-fin@cref}{{[ejercicio][6][2]2.6}{[1][7][]7}}
\newlabel{ej:conteos}{{3.3}{9}{}{ejemplo.3}{}}
\newlabel{ej:conteos@cref}{{[ejemplo][3][]3.3}{[1][9][]9}}
\newlabel{eje:3-laplace}{{3.1}{10}{}{ejercicio.3.1}{}}
\newlabel{eje:3-laplace@cref}{{[ejercicio][1][3]3.1}{[1][10][]10}}
\newlabel{eje:3-u1}{{3.2}{10}{}{ejercicio.3.2}{}}
\newlabel{eje:3-u1@cref}{{[ejercicio][2][3]3.2}{[1][10][]10}}
\newlabel{eje:3-interp}{{3.3}{10}{}{ejercicio.3.3}{}}
\newlabel{eje:3-interp@cref}{{[ejercicio][3][3]3.3}{[1][10][]10}}
\newlabel{eje:3-limite}{{3.4}{10}{}{ejercicio.3.4}{}}
\newlabel{eje:3-limite@cref}{{[ejercicio][4][3]3.4}{[1][10][]10}}
\newlabel{eje:4-orden}{{4.1}{13}{}{ejercicio.4.1}{}}
\newlabel{eje:4-orden@cref}{{[ejercicio][1][4]4.1}{[1][13][]13}}
\newlabel{eje:4-softmax}{{4.2}{13}{}{ejercicio.4.2}{}}
\newlabel{eje:4-softmax@cref}{{[ejercicio][2][4]4.2}{[1][13][]13}}
\newlabel{eje:4-muestreo}{{4.3}{13}{}{ejercicio.4.3}{}}
\newlabel{eje:4-muestreo@cref}{{[ejercicio][3][4]4.3}{[1][13][]13}}
\newlabel{eje:4-potencia}{{4.4}{13}{}{ejercicio.4.4}{}}
\newlabel{eje:4-potencia@cref}{{[ejercicio][4][4]4.4}{[1][13][]13}}
\newlabel{eje:4-limites}{{4.5}{13}{}{ejercicio.4.5}{}}
\newlabel{eje:4-limites@cref}{{[ejercicio][5][4]4.5}{[1][13][]13}}
\newlabel{eje:4-lse}{{4.6}{13}{}{ejercicio.4.6}{}}
\newlabel{eje:4-lse@cref}{{[ejercicio][6][4]4.6}{[1][13][]13}}
\newlabel{ej:ppl}{{5.3}{14}{}{ejemplo.3}{}}
\newlabel{ej:ppl@cref}{{[ejemplo][3][]5.3}{[1][14][]14}}
\newlabel{eje:5-ppl1}{{5.1}{16}{}{ejercicio.5.1}{}}
\newlabel{eje:5-ppl1@cref}{{[ejercicio][1][5]5.1}{[1][15][]16}}
\newlabel{eje:5-perdida}{{5.2}{16}{}{ejercicio.5.2}{}}
\newlabel{eje:5-perdida@cref}{{[ejercicio][2][5]5.2}{[1][16][]16}}
\newlabel{eje:5-digitos}{{5.3}{16}{}{ejercicio.5.3}{}}
\newlabel{eje:5-digitos@cref}{{[ejercicio][3][5]5.3}{[1][16][]16}}
\newlabel{eje:5-ppl}{{5.4}{16}{}{ejercicio.5.4}{}}
\newlabel{eje:5-ppl@cref}{{[ejercicio][4][5]5.4}{[1][16][]16}}
\newlabel{eje:5-lagrange}{{5.5}{16}{}{ejercicio.5.5}{}}
\newlabel{eje:5-lagrange@cref}{{[ejercicio][5][5]5.5}{[1][16][]16}}
\newlabel{ap:soluciones}{{A}{17}{Soluciones de los ejercicios}{appendix.A}{}}
\newlabel{ap:soluciones@cref}{{[appendix][1][2147483647]A}{[1][17][]17}}
\newlabel{sol:eje:1-contar}{{A.1}{17}{Tema 1}{section*.3}{}}
\newlabel{sol:eje:1-contar@cref}{{[subappendix][1][2147483647,1]A.1}{[1][17][]17}}
\newlabel{sol:eje:1-inyectiva}{{A.1}{17}{Tema 1}{section*.4}{}}
\newlabel{sol:eje:1-inyectiva@cref}{{[subappendix][1][2147483647,1]A.1}{[1][17][]17}}
\newlabel{sol:eje:1-granularidad}{{A.1}{17}{Tema 1}{section*.5}{}}
\newlabel{sol:eje:1-granularidad@cref}{{[subappendix][1][2147483647,1]A.1}{[1][17][]17}}
\newlabel{sol:eje:2-casillas}{{A.2}{18}{Tema 2}{section*.6}{}}
\newlabel{sol:eje:2-casillas@cref}{{[subappendix][2][2147483647,1]A.2}{[1][18][]18}}
\newlabel{sol:eje:2-vf}{{A.2}{18}{Tema 2}{section*.7}{}}
\newlabel{sol:eje:2-vf@cref}{{[subappendix][2][2147483647,1]A.2}{[1][18][]18}}
\newlabel{sol:eje:2-sol-matriz}{{A.2}{18}{Tema 2}{section*.8}{}}
\newlabel{sol:eje:2-sol-matriz@cref}{{[subappendix][2][2147483647,1]A.2}{[1][18][]18}}
\newlabel{sol:eje:2-sol-prob}{{A.2}{19}{Tema 2}{section*.9}{}}
\newlabel{sol:eje:2-sol-prob@cref}{{[subappendix][2][2147483647,1]A.2}{[1][19][]19}}
\newlabel{sol:eje:2-longitud}{{A.2}{19}{Tema 2}{section*.10}{}}
\newlabel{sol:eje:2-longitud@cref}{{[subappendix][2][2147483647,1]A.2}{[1][19][]19}}
\newlabel{sol:eje:2-sin-fin}{{A.2}{20}{Tema 2}{section*.11}{}}
\newlabel{sol:eje:2-sin-fin@cref}{{[subappendix][2][2147483647,1]A.2}{[1][20][]20}}
\newlabel{sol:eje:3-laplace}{{A.3}{20}{Tema 3}{section*.12}{}}
\newlabel{sol:eje:3-laplace@cref}{{[subappendix][3][2147483647,1]A.3}{[1][20][]20}}
\newlabel{sol:eje:3-u1}{{A.3}{20}{Tema 3}{section*.13}{}}
\newlabel{sol:eje:3-u1@cref}{{[subappendix][3][2147483647,1]A.3}{[1][20][]20}}
\newlabel{sol:eje:3-interp}{{A.3}{21}{Tema 3}{section*.14}{}}
\newlabel{sol:eje:3-interp@cref}{{[subappendix][3][2147483647,1]A.3}{[1][21][]21}}
\newlabel{sol:eje:3-limite}{{A.3}{21}{Tema 3}{section*.15}{}}
\newlabel{sol:eje:3-limite@cref}{{[subappendix][3][2147483647,1]A.3}{[1][21][]21}}
\newlabel{sol:eje:4-orden}{{A.4}{22}{Tema 4}{section*.16}{}}
\newlabel{sol:eje:4-orden@cref}{{[subappendix][4][2147483647,1]A.4}{[1][21][]22}}
\newlabel{sol:eje:4-softmax}{{A.4}{22}{Tema 4}{section*.17}{}}
\newlabel{sol:eje:4-softmax@cref}{{[subappendix][4][2147483647,1]A.4}{[1][22][]22}}
\newlabel{sol:eje:4-muestreo}{{A.4}{22}{Tema 4}{section*.18}{}}
\newlabel{sol:eje:4-muestreo@cref}{{[subappendix][4][2147483647,1]A.4}{[1][22][]22}}
\newlabel{sol:eje:4-potencia}{{A.4}{22}{Tema 4}{section*.19}{}}
\newlabel{sol:eje:4-potencia@cref}{{[subappendix][4][2147483647,1]A.4}{[1][22][]22}}
\newlabel{sol:eje:4-limites}{{A.4}{23}{Tema 4}{section*.20}{}}
\newlabel{sol:eje:4-limites@cref}{{[subappendix][4][2147483647,1]A.4}{[1][23][]23}}
\newlabel{sol:eje:4-lse}{{A.4}{23}{Tema 4}{section*.21}{}}
\newlabel{sol:eje:4-lse@cref}{{[subappendix][4][2147483647,1]A.4}{[1][23][]23}}
\newlabel{sol:eje:5-ppl1}{{A.5}{24}{Tema 5}{section*.22}{}}
\newlabel{sol:eje:5-ppl1@cref}{{[subappendix][5][2147483647,1]A.5}{[1][24][]24}}
\newlabel{sol:eje:5-perdida}{{A.5}{24}{Tema 5}{section*.23}{}}
\newlabel{sol:eje:5-perdida@cref}{{[subappendix][5][2147483647,1]A.5}{[1][24][]24}}
\newlabel{sol:eje:5-digitos}{{A.5}{24}{Tema 5}{section*.24}{}}
\newlabel{sol:eje:5-digitos@cref}{{[subappendix][5][2147483647,1]A.5}{[1][24][]24}}
\newlabel{sol:eje:5-ppl}{{A.5}{24}{Tema 5}{section*.25}{}}
\newlabel{sol:eje:5-ppl@cref}{{[subappendix][5][2147483647,1]A.5}{[1][24][]24}}
\newlabel{sol:eje:5-lagrange}{{A.5}{25}{Tema 5}{section*.26}{}}
\newlabel{sol:eje:5-lagrange@cref}{{[subappendix][5][2147483647,1]A.5}{[1][25][]25}}
\makeatother

\begin{document}
\frontmatter
\begin{titlepage}
\centering
\vspace*{2.5cm}
{\large\scshape Anatomía de un LLM}\par
\vspace{0.4cm}
{\normalsize Fundamentos matemáticos de los modelos de lenguaje}\par
\vspace{2.2cm}
\rule{\textwidth}{0.8pt}\par\vspace{0.5cm}
{\small\scshape Fase I \textperiodcentered\ Fundamentos matemáticos}\par\vspace{0.3cm}
{\Huge\bfseries Módulo 1\par}\vspace{0.3cm}
{\LARGE Modelos de lenguaje probabilísticos\par}
\vspace{0.5cm}\rule{\textwidth}{0.4pt}\par
\vspace{1.6cm}
{\large Resumen del módulo}\par
\vfill
{\small Edición de octubre de 2026}\par
\end{titlepage}

\chapter*{Sobre este resumen}
\addcontentsline{toc}{chapter}{Sobre este resumen}
\markboth{Sobre este resumen}{}

Este documento reúne lo indispensable de los apuntes del Módulo 1: el texto expositivo, todas las definiciones y todos los resultados (teoremas, proposiciones, lemas y corolarios), las tablas y las observaciones de contenido. Se han suprimido las demostraciones, los ejemplos, las lecturas, las figuras, los ejercicios y sus soluciones. La numeración de temas, secciones, enunciados y ecuaciones es la misma que en los apuntes, de modo que cualquier referencia puede buscarse allí; las que remiten a elementos suprimidos (un ejemplo, una figura, una ecuación de una demostración) conservan su número en los apuntes.

\section*{Notación}
\begin{center}
\begin{tabular}{@{}ll@{}}
\toprule
Símbolo & Significado \\
\midrule
$\V$, $K=\abs{\V}$ & vocabulario y su tamaño \\
$\V^{*}$ & conjunto de secuencias finitas de elementos de $\V$ \\
$x_{1:T}=(x_1,\dots,x_T)$ & secuencia de $T$ tokens \\
\fin & símbolo de inicio y fin de secuencia \\
$N(u,v)$ & número de apariciones del bigrama $(u,v)$ en el corpus \\
$z\in\R^{K}$, $T>0$ & vector de logits y temperatura \\
$H(p)$, $H(p,q)$ & entropía y entropía cruzada \\
$\ln$, $\log_2$ & logaritmo natural (nats) y en base 2 (bits) \\
\bottomrule
\end{tabular}
\end{center}


\tableofcontents

\mainmatter
% =====================================================================
\chapter{Vocabulario, tokenización y modelos de lenguaje}
\label{tema:1}
% =====================================================================

Este tema fija los objetos sobre los que trabaja todo el curso: el vocabulario, las secuencias de tokens y la distribución de probabilidad que llamamos modelo de lenguaje. La presentación sigue la introducción del capítulo~3 de~\cite{jm3} y la clase~4 de~\cite{cs224n}.

\section{Vocabulario y secuencias}

\begin{definicion}[Vocabulario]
Un \emph{vocabulario} es un conjunto finito no vacío $\V$ cuyos elementos se llaman \emph{tokens}. Se escribe $K=\abs{\V}$ y se fija una biyección $\V\to\{0,1,\dots,K-1\}$ que asigna a cada token su \emph{índice}.
\end{definicion}

\begin{definicion}[Secuencias]
Para $T\ge 0$, $\V^{T}$ denota el conjunto de secuencias $x_{1:T}=(x_1,\dots,x_T)$ con $x_t\in\V$; $\V^{0}$ contiene solo la secuencia vacía. El conjunto de todas las secuencias finitas es
\[
\V^{*}=\bigcup_{T\ge 0}\V^{T}.
\]
\end{definicion}

\begin{proposicion}\label{prop:contar-secuencias}
Si $K=\abs{\V}\ge 2$, el número de secuencias de longitud a lo sumo $L$ es
\[
\sum_{T=0}^{L} K^{T}=\frac{K^{L+1}-1}{K-1}.
\]
En particular, $\V^{*}$ es infinito numerable.
\end{proposicion}

\section{Tokenización}

El texto se representa como una cadena de caracteres de un alfabeto $\Sigma$ (por ejemplo, Unicode). Un modelo de lenguaje no opera sobre $\Sigma^{*}$, sino sobre secuencias de tokens.

\begin{definicion}[Tokenizador]
Un \emph{tokenizador} es una aplicación $\tau\colon\Sigma^{*}\to\V^{*}$ junto con una aplicación de \emph{decodificación} $\delta\colon\V^{*}\to\Sigma^{*}$ tal que $\delta(\tau(s))=s$ para toda cadena $s$. La condición implica que $\tau$ es inyectiva: dos textos distintos nunca producen la misma secuencia de tokens.
\end{definicion}

\addtocounter{definicion}{1}%

La elección de $\V$ determina la granularidad. El \cref{tab:granularidad} resume las tres opciones habituales; los modelos actuales usan subpalabras obtenidas con el algoritmo BPE, que se estudia en el módulo~12.

\begin{table}[h]
\centering\small
\caption{Granularidad de la tokenización.}\label{tab:granularidad}
\begin{tabular}{@{}lccl@{}}
\toprule
Unidad & $K$ típico & Longitud de la secuencia & Limitación principal \\
\midrule
Carácter & $\sim 10^{2}$ & muy larga & secuencias costosas de procesar \\
Palabra & $>5\cdot 10^{5}$ & corta & palabras fuera del vocabulario \\
Subpalabra (BPE) & $3\cdot 10^{4}$–$2{,}6\cdot 10^{5}$ & intermedia & requiere entrenar el tokenizador \\
\bottomrule
\end{tabular}
\end{table}

\section{Modelos de lenguaje}

\begin{definicion}[Modelo de lenguaje]\label{def:ml}
Un \emph{modelo de lenguaje} sobre $\V$ es una distribución de probabilidad $P$ sobre $\V^{*}$, es decir, una función $P\colon\V^{*}\to[0,1]$ tal que
\[
\sum_{x\in\V^{*}}P(x)=1.
\]
\end{definicion}

La suma es una serie, porque $\V^{*}$ es infinito (\cref{prop:contar-secuencias}). Que una construcción concreta cumpla la condición de normalización no es automático: el \cref{tema:2} muestra que un modelo de bigramas sin símbolo de fin la incumple.

\begin{observacion}
En la práctica no se trabaja con $P(x)$ directamente, sino con las probabilidades condicionadas del token siguiente, $P(x_t\mid x_{1:t-1})$. Un gran modelo de lenguaje es una función con parámetros $\theta$ que recibe el contexto $x_{1:t-1}$ y devuelve un vector de $K$ números reales (los \emph{logits}); la función softmax (\cref{tema:4}) los convierte en una distribución sobre $\V$. La regla de la cadena (\cref{tema:2}) conecta ambas descripciones.
\end{observacion}


% =====================================================================
\chapter{Regla de la cadena y modelos de Markov}
\label{tema:2}
% =====================================================================

Este tema descompone la probabilidad de una secuencia en factores de la forma $P(x_t\mid x_{1:t-1})$ y presenta la aproximación de Markov que da lugar a los modelos de $n$-gramas. Los resultados de probabilidad siguen a~\cite[§1.3]{bt}; la parte de $n$-gramas, a~\cite[§3.1]{jm3}.

\section{Probabilidad condicionada}

\begin{definicion}[Probabilidad condicionada]
Sea $(\Omega,\mathcal{F},P)$ un espacio de probabilidad y sean $A,B\in\mathcal{F}$ con $P(A)>0$. La \emph{probabilidad de $B$ condicionada a $A$} es
\[
P(B\mid A)=\frac{P(A\cap B)}{P(A)}.
\]
\end{definicion}

Para $A$ fijo, la aplicación $B\mapsto P(B\mid A)$ es a su vez una medida de probabilidad sobre $(\Omega,\mathcal{F})$~\cite[§1.3]{bt}. Despejando de la definición se obtiene la \emph{regla del producto}, $P(A\cap B)=P(A)\,P(B\mid A)$.

\section{La regla de la cadena}

\begin{teorema}[Regla de la cadena]\label{teo:cadena}
Sean $A_1,\dots,A_T$ sucesos con $P(A_1\cap\dots\cap A_{T-1})>0$. Entonces
\[
P\Bigl(\,\bigcap_{t=1}^{T}A_t\Bigr)=\prod_{t=1}^{T}P\bigl(A_t\mid A_1\cap\dots\cap A_{t-1}\bigr),
\]
donde el factor con $t=1$ se interpreta como $P(A_1)$.
\end{teorema}

\begin{corolario}\label{cor:cadena-secuencias}
Para un modelo de lenguaje y una secuencia $x_{1:T}$ con $P(x_{1:T-1})>0$,
\[
P(x_{1:T})=\prod_{t=1}^{T}P(x_t\mid x_{1:t-1}).
\]
\end{corolario}

La igualdad del \cref{cor:cadena-secuencias} es exacta. Todo modelo de lenguaje autorregresivo, desde un modelo de bigramas hasta un Transformer, consiste en una forma concreta de calcular cada factor $P(x_t\mid x_{1:t-1})$.

\section{Hipótesis de Markov y modelos de \texorpdfstring{$n$}{n}-gramas}

El contexto $x_{1:t-1}$ crece con $t$, y el número de contextos posibles crece exponencialmente (\cref{prop:contar-secuencias}). La solución clásica consiste en truncar el contexto~\cite[§3.1.1]{jm3}.

\begin{definicion}[Modelo de $n$-gramas]\label{def:ngramas}
Sea $n\ge 1$. Un modelo de lenguaje satisface la \emph{hipótesis de Markov de orden $n-1$} si
\[
P(x_t\mid x_{1:t-1})=P(x_t\mid x_{t-n+1:t-1})\qquad\text{para todo } t,
\]
es decir, si la distribución del token siguiente depende solo de los $n-1$ tokens anteriores. El modelo se llama de \emph{unigramas} ($n=1$), de \emph{bigramas} ($n=2$) o de \emph{trigramas} ($n=3$).
\end{definicion}

\begin{definicion}[Matriz estocástica]
Una matriz $P\in\R^{K\times K}$ es \emph{estocástica por filas} si $P_{ij}\ge 0$ para todo $i,j$ y $\sum_{j}P_{ij}=1$ para todo $i$. Un modelo de bigramas queda determinado por la matriz estocástica $P_{ij}=P(x_t=j\mid x_{t-1}=i)$, llamada \emph{matriz de transición}.
\end{definicion}

\begin{proposicion}\label{prop:parametros}
Un modelo de $n$-gramas sobre un vocabulario de tamaño $K$ se especifica con una tabla de $K^{n}$ probabilidades, de las cuales $K^{n-1}(K-1)$ son libres.
\end{proposicion}

Con $K=5\cdot 10^{4}$ y $n=3$, la tabla tendría $1{,}25\cdot 10^{14}$ entradas. Esta explosión combinatoria es la razón por la que los modelos modernos sustituyen la tabla por una función con parámetros compartidos (módulos~9 y~13).

\addtocounter{definicion}{1}%

\addtocounter{figure}{1}%

\section{El símbolo de fin y la normalización}

La \cref{def:ml} exige que las probabilidades de \emph{todas} las secuencias finitas sumen~1. El símbolo \fin\ de fin de secuencia es imprescindible para ello~\cite[§3.1, ejercicio 3.5]{jm3}.

\begin{proposicion}\label{prop:longitud-fija}
Sea un modelo de bigramas sin símbolo de fin, con distribución inicial $\pi$ sobre $\V$ y matriz de transición estocástica $P$. Para cada longitud $L\ge 1$,
\[
\sum_{x_{1:L}\in\V^{L}}\pi(x_1)\prod_{t=2}^{L}P_{x_{t-1}x_t}=1.
\]
En consecuencia, la suma sobre todas las longitudes diverge y el modelo no define una distribución sobre $\V^{*}$.
\end{proposicion}

\begin{observacion}
Con el símbolo de fin, la generación se detiene la primera vez que aparece \fin, y cada longitud recibe solo una parte de la masa total. Si desde cualquier token se alcanza \fin\ con probabilidad positiva, la masa que queda en secuencias de longitud $L$ decrece geométricamente con $L$ y la suma total es~1; el argumento completo es el de las cadenas de Markov absorbentes~\cite[cap.~7]{bt}.
\end{observacion}


% =====================================================================
\chapter{Estimación por máxima verosimilitud y suavizado}
\label{tema:3}
% =====================================================================

Los parámetros de un modelo de $n$-gramas se estiman a partir de un corpus. Este tema demuestra que el estimador de máxima verosimilitud es la frecuencia relativa, analiza su principal defecto (las probabilidades nulas) y presenta dos remedios: el suavizado aditivo y la interpolación. Se sigue~\cite[§3.1.2 y §3.6]{jm3} y, para la interpretación bayesiana, \cite[§2.2]{bishop}.

\section{Verosimilitud de un corpus}

Sea un corpus formado por secuencias, cada una delimitada por \fin\ al principio y al final, y sea $N(u,v)$ el número de veces que el token $v$ sigue inmediatamente al token $u$. Se escribe $N(u)=\sum_{v}N(u,v)$ para el número de veces que $u$ aparece como token anterior.

\begin{definicion}[Verosimilitud]
Dada una matriz de transición $P$, la \emph{verosimilitud} del corpus es el producto de las probabilidades de todos sus bigramas,
\[
\mathcal{L}(P)=\prod_{u,v\in\V}P_{uv}^{\,N(u,v)},
\qquad
\ell(P)=\ln\mathcal{L}(P)=\sum_{u,v\in\V}N(u,v)\ln P_{uv},
\]
y $\ell$ se llama \emph{log-verosimilitud}. El \emph{estimador de máxima verosimilitud} (EMV) es la matriz estocástica que maximiza $\ell$.
\end{definicion}

El logaritmo convierte productos en sumas y evita el subdesbordamiento numérico: con $10^{3}$ factores iguales a $0{,}9$ el producto vale aproximadamente $10^{-46}$.

\section{El estimador de máxima verosimilitud}

\begin{teorema}[EMV del modelo de bigramas]\label{teo:emv}
Supóngase $N(u)>0$ para todo $u$. El estimador de máxima verosimilitud de la matriz de transición es la frecuencia relativa
\[
\widehat{P}_{uv}=\frac{N(u,v)}{N(u)}.
\]
\end{teorema}

\addtocounter{definicion}{1}%

\begin{table}[h]
\centering\small
\caption{Conteos $N(u,v)$ del \cref{ej:conteos} (fila: token anterior; columna: token siguiente).}\label{tab:conteos}
\begin{tabular}{@{}c|ccccc|c@{}}
\toprule
$u\backslash v$ & \fin & \texttt{a} & \texttt{l} & \texttt{n} & \texttt{u} & $N(u)$ \\
\midrule
\fin & 0 & 2 & 1 & 0 & 0 & 3 \\
\texttt{a} & 3 & 0 & 1 & 1 & 0 & 5 \\
\texttt{l} & 0 & 1 & 0 & 0 & 1 & 2 \\
\texttt{n} & 0 & 2 & 0 & 0 & 0 & 2 \\
\texttt{u} & 0 & 0 & 0 & 1 & 0 & 1 \\
\bottomrule
\end{tabular}
\end{table}

\section{Ceros y suavizado aditivo}

El EMV asigna probabilidad nula a todo bigrama que no aparece en el corpus. Si un texto de evaluación contiene uno de ellos, su probabilidad es~0 y la pérdida del \cref{tema:5} es infinita. Este fenómeno es una forma de sobreajuste~\cite[§3.5]{jm3}.

\begin{definicion}[Suavizado aditivo]\label{def:laplace}
Sea $\alpha>0$. El estimador con \emph{suavizado aditivo} es
\[
\widehat{P}_{\alpha}(v\mid u)=\frac{N(u,v)+\alpha}{N(u)+\alpha K}.
\]
Con $\alpha=1$ recibe el nombre de \emph{suavizado de Laplace}; con $0<\alpha<1$, de \emph{suavizado add-$k$}~\cite[§3.6.1–3.6.2]{jm3}.
\end{definicion}

\begin{proposicion}\label{prop:laplace}
Para cada $u$, $\widehat{P}_{\alpha}(\cdot\mid u)$ es una distribución de probabilidad con todas sus entradas estrictamente positivas. Además, coincide con la media a posteriori de la fila $u$ cuando se adopta una distribución a priori de Dirichlet de parámetros $(\alpha,\dots,\alpha)$.
\end{proposicion}

\addtocounter{definicion}{1}%

\section{Interpolación}

Otra estrategia consiste en combinar estimadores de distinto orden~\cite[§3.6.3]{jm3}.

\begin{definicion}[Interpolación lineal]
Dados un estimador de bigramas $\widehat{P}_2(v\mid u)$ y uno de unigramas $\widehat{P}_1(v)$, y pesos $\lambda_1,\lambda_2\ge 0$ con $\lambda_1+\lambda_2=1$, el estimador \emph{interpolado} es
\[
\widehat{P}_{\mathrm{int}}(v\mid u)=\lambda_2\,\widehat{P}_2(v\mid u)+\lambda_1\,\widehat{P}_1(v).
\]
\end{definicion}

Es una distribución, por ser combinación convexa de distribuciones. Los pesos se eligen maximizando la verosimilitud de un conjunto de validación separado del de entrenamiento.


% =====================================================================
\chapter{Logits, softmax y generación}
\label{tema:4}
% =====================================================================

Una red neuronal no produce probabilidades directamente, sino un vector de números reales sin restricciones. Este tema estudia la función softmax, que transforma ese vector en una distribución, su versión con temperatura y los procedimientos de muestreo con los que un modelo genera texto. Se sigue~\cite[§4.1]{d2l}, \cite[§4.1]{gbc} y~\cite[§3.4]{jm3}.

\section{La función softmax}

\begin{definicion}[Softmax con temperatura]\label{def:softmax}
Sea $z\in\R^{K}$ un vector de \emph{logits} y $T>0$ una \emph{temperatura}. Se define
\[
\softmax_T(z)_i=\frac{e^{z_i/T}}{\sum_{j=1}^{K}e^{z_j/T}},\qquad i=1,\dots,K.
\]
Con $T=1$ se escribe simplemente $\softmax(z)$.
\end{definicion}

\begin{proposicion}[Propiedades]\label{prop:softmax}
Sean $z\in\R^{K}$ y $T>0$.
\begin{enumerate}
  \item $\softmax_T(z)$ pertenece al símplex: sus componentes son positivas y suman~1.
  \item Invariancia por traslación: $\softmax_T(z+c\mathbf{1})=\softmax_T(z)$ para todo $c\in\R$.
  \item Monotonía: si $z_i>z_j$, entonces $\softmax_T(z)_i>\softmax_T(z)_j$.
  \item Si el máximo $z_k$ es único, $\lim_{T\to 0^{+}}\softmax_T(z)=e_k$, el vector con un~1 en la posición $k$.
  \item $\lim_{T\to\infty}\softmax_T(z)=(1/K,\dots,1/K)$.
\end{enumerate}
\end{proposicion}

El \cref{tab:temperatura} ilustra la proposición para $z=(2,1,0)$. La temperatura no altera el orden de los tokens (propiedad~(iii)); solo controla cuánto se concentra la distribución.

\begin{table}[h]
\centering\small
\caption{$\softmax_T(z)$ para $z=(2,1,0)$.}\label{tab:temperatura}
\begin{tabular}{@{}lccc@{}}
\toprule
Temperatura & $p_1$ & $p_2$ & $p_3$ \\
\midrule
$T\to 0^{+}$ & $1$ & $0$ & $0$ \\
$T=0{,}5$ & $0{,}867$ & $0{,}117$ & $0{,}016$ \\
$T=1$ & $0{,}665$ & $0{,}245$ & $0{,}090$ \\
$T=2$ & $0{,}506$ & $0{,}307$ & $0{,}186$ \\
$T\to\infty$ & $1/3$ & $1/3$ & $1/3$ \\
\bottomrule
\end{tabular}
\end{table}

\begin{observacion}[Relación con los modelos de $n$-gramas]
Si en una fila de un modelo de bigramas se toman como logits $z_v=\ln\widehat{P}(v\mid u)$ para los $v$ con probabilidad positiva, entonces $\softmax(z)=\widehat{P}(\cdot\mid u)$. Elevar las probabilidades a $1/T$ y renormalizar equivale a aplicar la softmax con temperatura $T$.
\end{observacion}

\begin{observacion}[Relación con la física]
La expresión de la \cref{def:softmax} es la distribución de Boltzmann de la mecánica estadística, $p_i\propto e^{-E_i/(k_BT)}$, con energías $E_i=-z_i$ y la temperatura medida en unidades de $k_B$. Los límites (iv) y (v) corresponden al estado fundamental a temperatura nula y a la equiprobabilidad a temperatura infinita.
\end{observacion}

\section{Estabilidad numérica}

Para valores grandes de $z_i$, el cálculo directo de $e^{z_i}$ desborda la aritmética de coma flotante: en doble precisión el mayor número representable es aproximadamente $1{,}8\cdot 10^{308}$, y $e^{710}$ ya lo supera~\cite[§4.1]{gbc}.

\begin{definicion}[Log-suma-exp]
$\LSE(z)=\ln\sum_{i=1}^{K}e^{z_i}$. Se cumple $\ln\softmax(z)_i=z_i-\LSE(z)$.
\end{definicion}

\begin{proposicion}\label{prop:lse}
Sea $m=\max_i z_i$. Entonces $\LSE(z)=m+\LSE(z-m\mathbf{1})$ y
\[
m\le\LSE(z)\le m+\ln K.
\]
\end{proposicion}

Tras restar el máximo, todos los exponentes son $\le 0$ y no se produce desbordamiento. Las bibliotecas numéricas calculan así la softmax y la entropía cruzada.

\section{Muestreo}

Generar texto consiste en aplicar repetidamente el \cref{cor:cadena-secuencias}: se obtiene la distribución $p=P(\cdot\mid x_{1:t-1})$, se elige un token $x_t$ y se añade al contexto, hasta que aparece \fin~\cite[§3.4]{jm3}.

\begin{proposicion}[Muestreo por la inversa de la función de distribución]\label{prop:muestreo}
Sea $p=(p_1,\dots,p_K)$ una distribución, $F_k=\sum_{i\le k}p_i$ con $F_0=0$, y sea $U$ una variable aleatoria uniforme en $[0,1)$. Defínase $X=\min\{k: U<F_k\}$. Entonces $P(X=k)=p_k$ para todo $k$.
\end{proposicion}

\begin{definicion}[Estrategias de decodificación]
Dada $p=P(\cdot\mid x_{1:t-1})$:
\begin{enumerate}
  \item \emph{decodificación voraz}: $x_t=\arg\max_i p_i$ (equivale a $T\to 0^{+}$);
  \item \emph{muestreo top-$k$}: se conservan los $k$ tokens más probables, se renormaliza y se muestrea~\cite{fan};
  \item \emph{muestreo de núcleo} o \emph{top-$p$}: se conserva el menor conjunto de tokens, ordenados por probabilidad decreciente, cuya probabilidad acumulada alcanza un umbral $p_0\in(0,1]$, se renormaliza y se muestrea~\cite{holtzman}.
\end{enumerate}
\end{definicion}

Estas estrategias, combinadas con la temperatura, son los controles que exponen las interfaces de los modelos actuales. Se estudian con más detalle en el módulo~19.


% =====================================================================
\chapter{Evaluación: entropía cruzada y perplejidad}
\label{tema:5}
% =====================================================================

Un buen modelo asigna probabilidad alta al texto real. Este tema define la pérdida que se minimiza al entrenar todos los modelos de lenguaje, la entropía cruzada, y la medida de evaluación derivada de ella, la perplejidad. Se sigue~\cite[§3.2, §3.3 y §3.7]{jm3}; los resultados de teoría de la información proceden de~\cite[cap.~2]{ct}.

\section{Pérdida y perplejidad}

Sea un texto de evaluación con $N$ predicciones y sea $p_t=P(x_t\mid x_{1:t-1})\in(0,1]$ la probabilidad que el modelo asigna al token que realmente aparece en la posición $t$.

\begin{definicion}[Pérdida y perplejidad]\label{def:ppl}
La \emph{pérdida} (log-verosimilitud negativa media) y la \emph{perplejidad} del modelo sobre el texto son
\[
\mathcal{J}=-\frac{1}{N}\sum_{t=1}^{N}\ln p_t,\qquad \PPL=e^{\mathcal{J}}.
\]
Con logaritmo natural, $\mathcal{J}$ se mide en \emph{nats}; con $\log_2$, en \emph{bits}, y $\mathcal{J}_{\text{bits}}=\mathcal{J}_{\text{nats}}/\ln 2$.
\end{definicion}

\begin{proposicion}\label{prop:ppl}
\begin{enumerate}
  \item $\PPL=\bigl(\prod_{t=1}^{N}1/p_t\bigr)^{1/N}$, la media geométrica de los inversos de las probabilidades.
  \item $\PPL\ge 1$, con igualdad si y solo si $p_t=1$ para todo $t$.
  \item Un modelo uniforme sobre $K$ tokens tiene $\PPL=K$ sobre cualquier texto.
\end{enumerate}
\end{proposicion}

La perplejidad admite una interpretación como \emph{factor de ramificación medio}: un modelo con perplejidad $k$ está, en promedio, tan indeciso como si eligiera uniformemente entre $k$ tokens~\cite[§3.3.1]{jm3}.

\addtocounter{definicion}{1}%

\section{Entropía y entropía cruzada}

\begin{definicion}[Entropía y entropía cruzada]
Sean $p$ y $q$ distribuciones sobre un conjunto finito $\mathcal{X}$, con $q(x)>0$ siempre que $p(x)>0$. La \emph{entropía} de $p$ y la \emph{entropía cruzada} de $q$ respecto de $p$ son
\[
H(p)=-\sum_{x\in\mathcal{X}}p(x)\ln p(x),\qquad
H(p,q)=-\sum_{x\in\mathcal{X}}p(x)\ln q(x),
\]
con el convenio $0\ln 0=0$.
\end{definicion}

Si los tokens del texto se distribuyen según $p$ y el modelo es $q$, la pérdida $\mathcal{J}$ es un estimador de $H(p,q)$: es la media empírica de $-\ln q(X)$ con $X\sim p$.

\begin{teorema}[Desigualdad de Gibbs]\label{teo:gibbs}
$H(p,q)\ge H(p)$, con igualdad si y solo si $q=p$.
\end{teorema}

La diferencia $H(p,q)-H(p)=\sum_x p(x)\ln\bigl(p(x)/q(x)\bigr)$ es la \emph{divergencia de Kullback–Leibler} $D(p\,\|\,q)$, que se estudia en el módulo~6.

\begin{corolario}\label{cor:aprender}
Entre todos los modelos $q$, la entropía cruzada esperada $H(p,q)$ se minimiza exactamente en $q=p$, y su valor mínimo es $H(p)$.
\end{corolario}

El \cref{cor:aprender} justifica el criterio de entrenamiento de los modelos de lenguaje: minimizar la entropía cruzada media sobre un corpus empuja al modelo a reproducir la distribución real del token siguiente. Ningún modelo puede bajar de $H(p)$, la incertidumbre propia del lenguaje. La relación entre entropía y predicción se remonta a Shannon~\cite{shannon}.

\section{Entrenamiento y evaluación}

\begin{observacion}
La perplejidad debe medirse sobre un conjunto de evaluación disjunto del de entrenamiento. Evaluar sobre los datos de entrenamiento da una estimación optimista: el EMV del \cref{teo:emv} maximiza precisamente la verosimilitud de esos datos~\cite[§3.2]{jm3}. Además, la perplejidad solo es comparable entre modelos que usan el mismo vocabulario, porque depende de la tokenización.
\end{observacion}



\backmatter
% =====================================================================
\begin{thebibliography}{99}
\addcontentsline{toc}{chapter}{Bibliografía}
\markboth{Bibliografía}{}
\small

\bibitem{bt}
D.~P. Bertsekas y J.~N. Tsitsiklis.
\newblock \emph{Introduction to Probability}, 2.ª ed.
\newblock Athena Scientific, 2008.
\newblock Texto del curso 6.041 del MIT.

\bibitem{bishop}
C.~M. Bishop.
\newblock \emph{Pattern Recognition and Machine Learning}.
\newblock Springer, 2006.

\bibitem{chen-goodman}
S.~F. Chen y J.~Goodman.
\newblock An empirical study of smoothing techniques for language modeling.
\newblock \emph{Computer Speech \& Language}, 13(4):359--394, 1999.

\bibitem{ct}
T.~M. Cover y J.~A. Thomas.
\newblock \emph{Elements of Information Theory}, 2.ª ed.
\newblock Wiley, 2006.

\bibitem{cs224n}
D.~Yang y Y.~Choi.
\newblock CS224N: Natural Language Processing with Deep Learning.
\newblock Stanford University, invierno de 2026. Clase~4, «Language Models and RNNs».
\newblock \url{https://web.stanford.edu/class/cs224n/}

\bibitem{fan}
A.~Fan, M.~Lewis e Y.~Dauphin.
\newblock Hierarchical neural story generation.
\newblock En \emph{Proceedings of ACL}, 2018.
\newblock \url{https://arxiv.org/abs/1805.04833}

\bibitem{gbc}
I.~Goodfellow, Y.~Bengio y A.~Courville.
\newblock \emph{Deep Learning}.
\newblock MIT Press, 2016.
\newblock \url{https://www.deeplearningbook.org/}

\bibitem{holtzman}
A.~Holtzman, J.~Buys, L.~Du, M.~Forbes e Y.~Choi.
\newblock The curious case of neural text degeneration.
\newblock En \emph{International Conference on Learning Representations (ICLR)}, 2020.
\newblock \url{https://arxiv.org/abs/1904.09751}

\bibitem{jm3}
D.~Jurafsky y J.~H. Martin.
\newblock \emph{Speech and Language Processing}, 3.ª ed., borrador del 19 de agosto de 2026.
\newblock Cap.~3, «N-gram Language Models».
\newblock \url{https://web.stanford.edu/~jurafsky/slp3/}

\bibitem{mackay}
D.~J.~C. MacKay.
\newblock \emph{Information Theory, Inference, and Learning Algorithms}.
\newblock Cambridge University Press, 2003.
\newblock \url{http://www.inference.org.uk/itila/book.html}

\bibitem{mit6041}
J.~N. Tsitsiklis.
\newblock 6.041SC: Probabilistic Systems Analysis and Applied Probability.
\newblock Massachusetts Institute of Technology, MIT OpenCourseWare, otoño de 2013.
\newblock \url{https://ocw.mit.edu/courses/6-041sc-probabilistic-systems-analysis-and-applied-probability-fall-2013/}

\bibitem{shannon}
C.~E. Shannon.
\newblock A mathematical theory of communication.
\newblock \emph{Bell System Technical Journal}, 27:379--423 y 623--656, 1948.

\bibitem{d2l}
A.~Zhang, Z.~C. Lipton, M.~Li y A.~J. Smola.
\newblock \emph{Dive into Deep Learning}, versión 1.0.3.
\newblock Cambridge University Press, 2023.
\newblock \url{https://d2l.ai/}

\end{thebibliography}


\end{document}
